\chapter*{Preface}
\addcontentsline{toc}{chapter}{Preface}

Among the vast types of questions within graph theory, I think
vertex-partitioning problems are the ones that I like the most, for they combine
two of my favorite aspects in mathematics: the puzzle-like structure I am very
fond of, and the challenge to find efficient algorithmic procedures to find
partitions.

When I was first introduced to this kind of problems, I got immediately into it.
So much was my interest that, in several occasions during the process of writing
this text, I could not find myself at peace without ensuring that a claim was
indeed true.

After a year or so of living that on a weekly basis, I can confidently say that,
beyond the mathematical results obtained, I have grown to like the \textit{loop}
in which we tackle these problems: an initial computational search for
obstructions to a hereditary property, leading to a partial proof of a
characterization, that in turn yield more obstructions, and so on and so forth.

That being said, the particular problem we take a look at is polarity in graphs.
Nonetheless, rather than working directly onto it, such a task is that hard to
consider as an starting point, so we restrict ourselves to solve a smaller issue
regarding monopolarity. In an effort to generalize previous results (and to
improve some time-bounds), we consider the family of graphs free of the five
vertex path and its complement, since it being a selfcomplementary class allows
us to automatically solve another problem similar to monopolarity.

We begin the thesis with the usual mandatory chapter of introduction, in which
we establish our framework and give some previous results on the matter. Later,
we divide our study into three main chapters, considering if there are
underlying structures within our graphs that could ease the analysis to be done:
first we check on chordal graphs, then we restrict ourselves to graphs with an
induced cycle of $5$-vertices, and we finish with graphs with an induced cycle
of $4$-vertices.

Lastly, we gather all of our results in the conclusions and leave some lines of
possible further research. Mainly, the existence of a linear-time certifying
algorithm to decide monopolarity on $(P_5,\text{house})$-free graphs shines a
light on the possibility of there being a similar algorithm deciding polarity in
the same class of graphs.